The two formulae for the highest and the lowest set bit look alike, and they do not cost the
same.
CPython stores a non-negative integer as its digits in base $2^{30}$,
together with their number.
\codehl{bit\_length} reads that number and the highest digit, so it costs $O(1)$.
An operator such as $\bitAnd$, $\bitOr$, $\oplus$ or a shift builds a new integer digit by
digit, and so does the negation $b \mapsto -b$;
on an integer of $\bandWidth$ bits each costs $O(\bandWidth)$,
and so does \codehl{bit\_count}, which counts the set bits of every digit.
Integers are immutable, so an update in place such as \codehl{b |= bit} builds a new integer
too, and costs $O(\bandWidth)$ as well.

So on a set held as an integer of $\bandWidth$ bits,
the largest element costs $O(1)$ and the smallest costs $O(\bandWidth)$.
The asymmetry belongs to the unbounded integers of Python and not to the algorithm:
on a machine word, both are a single instruction of the processor.
Listing \ref{listing.bitCosts} measures both on two widths a hundred times apart;
the first time does not grow with the width and the second does.

\begin{lstlisting}[caption={The highest set bit against the lowest, on two widths.}, label=listing.bitCosts]
import timeit

repeats = 20_000
for width in (10**3, 10**5):
    b = 1 << width
    highest = timeit.timeit(lambda: b.bit_length(), number=repeats) / repeats
    lowest = timeit.timeit(lambda: b & -b, number=repeats) / repeats
    print(f"{width:>6} bits: {highest * 1e9:7.1f} ns against {lowest * 1e9:9.1f} ns")
\end{lstlisting}
